%! Polar Function Graphs — Unit 3, Lesson 3.14 — Student Packet Cover Sheet
\documentclass[10pt]{article}
\usepackage{precalculus-article}
\usepackage{precalculus-boxes}
\usepackage{ltablex}
\keepXColumns

\begin{document}

% Full-bleed navy banner
\begin{tikzpicture}[remember picture, overlay]
  \fill[navy] (current page.north west)
    rectangle ([yshift=-0.9in]current page.north east);
\end{tikzpicture}
\vspace{-0.8in}

\begin{center}
  {\color{white}\textbf{\LARGE AP Precalculus}}\\[4pt]
  {\color{skymid}\large\textbf{Unit 3: Trigonometric and Polar Functions}}\\[6pt]
  {\color{white}\textbf{\large Lesson 3.14 \quad Polar Function Graphs}}
\end{center}

\vspace{0.22in}
\namedateperiod
\vspace{0.18in}

\begin{learningtargetbox}
\small
\begin{enumerate}[leftmargin=1.5em, itemsep=5pt]
  \item I can read a polar function $r = f(\theta)$ from a table of input-output pairs
    and identify the maximum radius, the minimum radius, and the $\theta$ values where
    the curve passes through the pole.
    \hfill\textit{(LO 3.14.A, EK 3.14.A.1)}
  \item I can connect a rectangular graph of $r = f(\theta)$ to the polar curve by
    identifying intervals where $r$ is positive or negative and using zeros and
    extrema to determine the number of petals in a rose curve.
    \hfill\textit{(LO 3.14.A, EK 3.14.A.3)}
  \item I can describe which portion of a polar curve is traced when the domain of
    $\theta$ is restricted to a sub-interval.
    \hfill\textit{(LO 3.14.A, EK 3.14.A.2)}
\end{enumerate}
\end{learningtargetbox}

\vspace{0.14in}

\begin{tocbox}
\small
\renewcommand{\arraystretch}{1.5}
\begin{tabularx}{\linewidth}{@{} c l X r @{}}
\rowcolor{sky}
\textbf{\#} & \textbf{Component} & \textbf{Description} & \textbf{Score} \\
\midrule
1 & Warm-Up        & Spiral review: unit-circle sine values, reading $r = 2\sin\theta$ from a rectangular graph, polar function input/output roles & \blank{1.2cm} \\
2 & Guided Notes   & Polar function vocabulary; reading from tables; connecting rectangular and polar graphs; domain restriction & \blank{1.2cm} \\
3 & Group Activity & Table reading (Tier R), rectangular graph interpretation (Tier A), comparing two polar curves (Tier E) & \blank{1.2cm} \\
4 & Exit Ticket    & Complete table for $r = 3\cos\theta$; identify zeros; MC on input/output roles & \blank{1.2cm} \\
5 & Homework       & Table questions, petal-count from rectangular graph, domain restriction, AP-style MC & \blank{1.2cm} \\
\midrule
  & \multicolumn{2}{r}{\textbf{Total}} & \blank{1.2cm} \\
\end{tabularx}
\end{tocbox}

\end{document}
